\chapter{Вырожденный символ директорного сектора и предел факторизации}

Предыдущая глава установила положительный $J=1$-порог с законом
$L^{-4}$. До попытки факторизовать все канальные операторы необходимо
вычислить их главный символ на однородном упорядоченном вакууме. Именно он
решает, является ли нулевой порог обычной дисперсией или следствием более
сильной вырожденности.

\section{Литературная и проектная граница}

Обычная модель Скирма содержит и квадратичный сигма-модельный, и
четырёхпроизводный скирмовский члены; их баланс отвечает за конечный размер
и жёсткость солитона. Современное обсуждение этой устойчивости дано в
\path{arXiv:2405.05731}. Сильно связанный предел Фаддеева--Скирма без
квадратичного кинетического члена является отдельной вырожденной моделью;
см. \path{arXiv:0911.3673}. Поэтому отсутствие отрицательной моды ещё не
означает, что чисто кривизный предел обладает невырожденной линейной
динамикой.

В текущем родителе двухпроизводный член формально присутствует как
$|D_QQ|^2$, но составная связность может уничтожать его на касательных
направлениях вакуумной орбиты. Это проверяется ниже без аналогий с готовыми
моделями.

\section{Символ гессиана на однородном вакууме}

Выберем
\begin{equation}
 Q_0=\Delta\left(P_0-\frac{I_3}{3}\right),
 \qquad P_0=\operatorname{diag}(1,0,0).
\end{equation}
Пять компонент $\delta Q$ распадаются на три нормальных к вакуумной
орбите и две касательные директорные компоненты.

Для плоской волны с волновым числом $k$ линейная вариация составной
связности равна
\begin{equation}
 \delta A_i=\frac{[Q_0,\partial_i\delta Q]}{\Delta^2}.
\end{equation}
Если $\delta Q$ касательна к $SO(3)/O(2)$, проекторное тождество даёт
\begin{equation}
 \delta(D_iQ)=
 \partial_i\delta Q+[\delta A_i,Q_0]=0.
 \label{eq:v6-tangent-DQ-symbol-zero}
\end{equation}
Поскольку фон постоянен и $A_i(Q_0)=0$, также
\begin{equation}
 \delta F_{ij}=\partial_i\delta A_j-partial_j\delta A_i=0.
 \label{eq:v6-tangent-F-symbol-zero}
\end{equation}
Локальный потенциал имеет две голдстоуновские нулевые кривизны. Поэтому
полный статический символ точно обращается в нуль на двух директорных
компонентах при любом $k$.

Машинное вычисление при $k=0,1/2,1,3$ дало ранг три из пяти. Две
номинально нулевые величины равны $2.28\cdot10^{-7}$ и полностью
объясняются шагом численного дифференцирования канонического потенциала;
точное орбитальное тождество делает их нулевыми.

\begin{theorem}[Неэллиптичность статического гессиана]
Главный символ канонического статического гессиана имеет точное
двумерное ядро директорных направлений для каждого ненулевого волнового
вектора. Поэтому полный пятикомпонентный оператор не является
эллиптическим.
\end{theorem}

\section{Точные плоские семейства, а не только нулевой гессиан}

Вырожденность сохраняется за пределами второй вариации. Пусть директор
$n=n(x^1)$ произвольно зависит только от одной координаты и
\begin{equation}
 Q(x)=\Delta\left(n(x^1)n(x^1)^T-\frac{I_3}{3}\right).
\end{equation}
Тогда
\begin{equation}
 D_QQ=0,
 \qquad V_{\rm can}(Q)=0.
\end{equation}
Связность имеет только одну пространственную компоненту, поэтому все
антисимметричные компоненты кривизны $F_{ij}$ также равны нулю. Получаем
\begin{equation}
 \mathcal E[Q]=0
 \label{eq:v6-exact-flat-director-family}
\end{equation}
для произвольной одномерной директорной волны.

Численная проверка на $41$ точке дала максимальный остаток
$|D_QQ|=2.2\cdot10^{-16}$ и точно нулевую кривизну. Для директора,
зависящего от двух координат, ковариантная производная всё ещё исчезает,
но кривизна уже ненулевая: в контрольной точке
$|F_{12}|=0.13447266\ldots$.

\begin{proposition}[Плоская слоистая директорная ветвь]
Текущий составной функционал допускает бесконечномерное семейство
непостоянных одномерных вакуумных конфигураций с точно нулевой статической
энергией. Кривизный член различает лишь изменения директора по двум и более
независимым направлениям.
\end{proposition}

\section{Уточнение закона четвёртой степени}

Закон $\lambda\sim L^{-4}$ предыдущей главы воспроизводится, но теперь его
область интерпретации стала точнее. Он не является свободной
бигармонической дисперсией директорного поля: свободный квадратичный символ
директорного сектора тождественно равен нулю. Ненулевые конечнобъёмные
значения возникают из взаимодействия плоского сектора с убывающей
кривизной самого ежа и с граничным обрезанием пробных функций.

Это не возвращает отрицательную моду. Оно показывает, что положительные
конечнобъёмные числа не сходятся к невырожденному оператору частичных
колебаний.

\section{Почему требуемая факторизация не замыкается}

Коэрцитивная факторизация должна была бы иметь вид
\begin{equation}
 \mathcal H_J=\mathcal B_J^*\mathcal B_J+\mathcal M_J,
 \qquad \mathcal M_J\ge m_J^2>0
\end{equation}
вне конечного числа симметрийных нулевых мод. Точное двумерное ядро
главного символа при каждом $k$ запрещает такую оценку. Проблема состоит не
в скрытой отрицательной кривизне, а в отсутствии квадратичной жёсткости
целого полевого сектора.

Полуопределённая факторизация формально не исключена, но она не сделает
дефект изолированным локальным минимумом и не даст фредгольмов статический
гессиан. Поэтому дальнейшее увеличение конечных базисов не способно
закрыть требуемое утверждение.

\begin{nogocriterion}[Строгая устойчивость текущего родителя закрыта]
Пока директорный главный символ имеет ядро для каждого волнового числа,
нельзя называть составной ёж строго линейно устойчивой локализованной
частицей. Отсутствие отрицательных мод означает только отсутствие
экспоненциального спуска в проверенных направлениях, но не наличие
восстанавливающей силы.
\end{nogocriterion}

\section{Архитектурное следствие}

Остаются два принципиальных ремонта.
\begin{enumerate}
 \item Родительский квадратичный член обычной директорной жёсткости сделал
 бы гессиан эллиптическим, но вернул бы инфракрасную расходимость глобального
 ежа, если вращение директора не калибровать.
 \item Независимая полная связность могла бы превратить локальные
 директорные изменения в настоящую калибровочную избыточность. Но для этого
 требуется ранее отсутствовавшее стабилизаторное направление, собственная
 кинетика связности и физическая калибровочная факторизация.
\end{enumerate}

Следующий гейт должен проверить второй путь как минимальное переоткрытие, а
не просто добавить произвольное поле.

\begin{protocol}
\begin{enumerate}
 \item ранг пятикомпонентного символа: \textbf{$3$};
 \item размерность директорного ядра: \textbf{$2$ при каждом $k$};
 \item эллиптичность статического гессиана: \textbf{нет};
 \item одномерные директорные волны: \textbf{точно плоские};
 \item двумерные директорные текстуры: \textbf{видят кривизну};
 \item прежний закон $L^{-4}$: \textbf{сохраняется};
 \item свободная бигармоническая дисперсия: \textbf{не выведена};
 \item отрицательная мода: \textbf{не найдена};
 \item коэрцитивная канальная факторизация: \textbf{запрещена};
 \item строгая линейная устойчивость частицы: \textbf{не получена};
 \item причина провала: \textbf{плоскость, а не отрицательная кривизна};
 \item следующий гейт: \textbf{минимальное полное калибровочное
 завершение}.
\end{enumerate}
\end{protocol}

Машинный сертификат сохранён в
\path{s2t_v6_bosonic_defect_channel_operator_factorization_gate_results.json}.